\begin{abstract}
  A rapid increase in continuous seismic data volumes demands automated 
  solutions that can operate efficiently at regional scales. This work 
  evaluates an end-to-end, scalable \ac{ML} pipeline for earthquake detection 
  and characterization, designed for deployment in a \ac{HPC} environment. The 
  pipeline integrates publicly available models from the seismological 
  community via the SeisBench framework (\cite{SeisBench}) to process waveform 
  data from the Northeast Italy Monitoring System (\cite{SMINO}), managed by 
  the \acf{OGS}, as well as data from other Italian and international seismic
  networks in neighboring regions and countries.

  The workflow includes specialized modules for phase picking (PhaseNet, 
  \cite{PhaseNet}; EQTransformer, \cite{EQTransformer}), event association 
  (GaMMA, \cite{GaMMA}; PyOcto, \cite{PyOcto}), hypocenter location (NonLinLoc,
  \cite{NonLinLoc}), and magnitude estimation. Deep-learning phase pickers for 
  Northeastern Italy are extensively evaluated using diverse training datasets, 
  including INSTANCE (\cite{INSTANCE}), STEAD (\cite{STEAD}), SCEDC 
  (\cite{SCEDC}), and the original PhaseNet and EQTransformer training dataset,
  to determine the most effective combinations and ensure robust detection
  across the region.

  Performance is rigorously validated against the OGS reference catalogs using
  strict spatiotemporal matching criteria and recall-based metrics.

  Results show that the pipeline achieves high fidelity in detecting P- and 
  S-wave arrivals and in associating events, substantially outperforming 
  traditional CPU-only implementations in both speed and scalability. The 
  PhaseNet model trained on the INSTANCE dataset achieves the best overall 
  performance. The HPC architecture provides enhanced operational flexibility,
  enabling rapid near-real-time analysis as well as large-scale historical
  reprocessing on the same platform.
\end{abstract}